\clearpage\setcounter{chapter}{3}
\chapter{Hilbert 复形}\label{chap:hilbert-complexes}

刚才我们已经看到, de Rham 复形中的微分算子可以视为 $L^2$ 空间之间的闭无界算子. 这样便得到一种结构, 它将 Chapter~\ref{chap:homological-algebra} 中讨论的上链复形的同调代数特征, 与 Chapter~\ref{chap:unbounded-operators} 中讨论的 Hilbert 空间上闭无界算子的泛函分析特征结合起来. 这一结构称为 \emph{Hilbert 复形}, 最早由 Brüning 和 Lesch \BookCite{24} 提出. 它是 Riemannian 流形上 Hodge 理论的基础结构,%
\footnote{{Hilbert 复形概念的自然性, 可以从下面这件事看出: 在 \BookCite{13} 的早期草稿中, 我和合作者引入了这个定义, 此后才发现文献 \BookCite{24} 已经用完全相同的名称和定义提出了这一概念.}}%
而且, 尽管这一结构简单, Hodge 理论的主要结论, 例如将上同调与调和形式等同以及 Hodge 分解, 都能在该框架中容易地得到. 每个 Hilbert 复形都伴随一个分次映射或算子序列, 称为该复形的抽象 Hodge Laplacian. 我们将考察 Hodge Laplacian 的几种不同但等价的表述, 并说明在相差调和形式的意义下它们是适定的. 三维区域上的 de Rham 复形情形将得到详细讨论, 本章结果概括于章末的 Box~\ref{box:hilbert-complexes-summary}.
